% Reorganized from the frozen October 15, 2026 supplement.
\section{Probability tools and reporting conventions}
\label{app:tools}

For the reporting constructions, index arms in arrival order
$i=1,2,\ldots$, with a fixed order within each batch.
Every arm has an independent potential reward stream $(Y_{i,n})_{n\ge1}$
in $[0,1]$. The policy observes its $n$th entry on its $n$th pull. Arrival
times are deterministic, and reports are issued after the current observation.
Write $N_i(t)$ for lifetime observations through round $t$ and
$\hat r_i(n)=n^{-1}\sum_{s\le n}Y_{i,s}$. The observed reward in physical-time notation is
$X_{t,I_t}=Y_{I_t,N_{I_t}(t-1)+1}$. The conclusions below also
hold conditionally on an independent exogenous arrival schedule.

For independent $[0,1]$ observations, Hoeffding's inequality gives
\[
 P\{|\hat r_i(n)-r_i|>x\}\le2e^{-2nx^2}.
\]
For completeness, the centered log moment generating function has value and
derivative zero at the origin; its second derivative is a variance under an
exponentially tilted distribution and is at most $1/4$. Integrating twice
gives $Ee^{\lambda(X-EX)}\le e^{\lambda^2/8}$. Multiplication over independent
observations and exponential Markov inequality prove the display.

With
\[
 b_i(n)=\sqrt{\frac{\log(8i^2n^2/\delta)}{2n}},
\]
the union of all potential-prefix failures has probability at most
$\delta\zeta(2)^2/4<\delta$. Intersecting the prefix intervals with $[0,1]$
therefore gives simultaneous coverage. On that event each endpoint differs
from its mean by at most $2b_i(n)$. No conditioning on a previously selected
empirical mean is used. Empty intersections on failure paths are handled by
exclusion-first reporting; the reward policy remains defined on all paths.

The checkpoint alternative fixes, before observing arm $i$,
\[
 q_i(\rho_i)=\left\lceil\frac2{\rho_i^2}
            \log\frac{4i(i+1)}\delta\right\rceil.
\]
Before $q_i$ observations its interval is $[0,1]$; afterwards it is permanently
\[
 [\hat r_i(q_i)-\rho_i/2,
  \hat r_i(q_i)+\rho_i/2]\cap[0,1].
\]
The sum of coverage failures is at most
$\sum_i\delta/\{2i(i+1)\}=\delta/2$. On coverage, each endpoint differs from
the true mean by at most $\rho_i$. These intervals are nonempty on every path
and only contract once. If $\underline r^*(t)=\max_i\ell_i(t)$ and
$\overline r^*(t)=\max_i u_i(t)$, use
\[
 \mathcal P_t=\{i:\ell_i(t)\ge\overline r^*(t)-\varepsilon\},\quad
 \mathcal N_t=\{i:u_i(t)<\underline r^*(t)-\varepsilon\},\quad
 \mathcal V_t=\mathcal{A}_t\setminus(\mathcal P_t\cup \mathcal N_t).
\]
The sets $\mathcal P_t,\mathcal N_t$ are disjoint because
$\underline r^*(t)\le\overline r^*(t)$ and $\ell_i\le u_i$.
Coverage implies
$\mathcal P_t\subseteq \mathcal G_t^\varepsilon\subseteq \mathcal{A}_t\setminus \mathcal N_t$.
Moreover, $\mathcal N_t$ is increasing as observations and arrivals accumulate:
$\underline r^*(t)$ cannot decrease and an existing $u_i$ cannot increase.

For two Bernoulli mean vectors $\boldsymbol r=(r_i)_i$ and
$\boldsymbol r'=(r_i')_i$ and a bounded stopping time $\tau$, action
kernels, policy randomization, and stopping kernels cancel in the history
likelihood ratio. Padding after stopping by null symbols and taking iterated
expectations yields
\begin{equation}
 D(P_{\boldsymbol r}^{H_\tau}\Vert P_{\boldsymbol r'}^{H_\tau})
 =\sum_iE_{\boldsymbol r} N_i(\tau)\operatorname{kl}(r_i,r_i').
 \label{eq:app_kl_identity}
\end{equation}
The same proof applies to arbitrary observation laws with finite two-way KL.
For an event $F$ measurable at stopping, splitting the likelihood integral
over $F,F^c$ and applying Jensen gives
\[
 D(P\Vert Q)\ge\operatorname{kl}(P(F),Q(F))
 \ge P(F)\log\frac1{Q(F)}-\log2.
\]
The second inequality uses the binary entropy bound $h(p)\le\log2$. We repeatedly use
\[
 2(p-q)^2\le\operatorname{kl}(p,q)
 \le\frac{(p-q)^2}{q(1-q)}.
\]
The upper bound follows from $\log x\le x-1$; for the lower bound the second
derivative in $p$ is $1/\{p(1-p)\}\ge4$ and both value and derivative vanish
at $p=q$.
